% Generated draft: requires mathematical and source review.
\documentclass[11pt]{article}
\usepackage[margin=25mm]{geometry}
\usepackage{fontspec}
\setmainfont{Latin Modern Roman}
\usepackage{amsmath,amssymb,mathtools,unicode-math}
\setmathfont{Latin Modern Math}
\usepackage{enumitem,needspace,xurl,hyperref}
\setlength{\parindent}{0pt}
\setlength{\parskip}{5pt}
\setlength{\emergencystretch}{3em}
\newcommand{\N}{\mathbb{N}}
\newcommand{\Z}{\mathbb{Z}}
\newcommand{\Q}{\mathbb{Q}}
\newcommand{\R}{\mathbb{R}}
\newcommand{\C}{\mathbb{C}}
\newcommand{\F}{\mathbb{F}}
\newcommand{\A}{\mathbb{A}}
\newcommand{\PP}{\mathbb{P}}
\newcommand{\E}{\mathbb{E}}
\newcommand{\eps}{\varepsilon}
\newcommand{\dd}{\,\mathrm{d}}
\begin{document}
\section*{UnsolvedMath ID: 30004553}
\textbf{Source problem code:} catalogue entry
\subsection{Definitions and mathematical statement}
Let \(\mathrm{MF}\) denote the Minimalist Foundation referred to in the supplied source record, and let \(\mathrm{MF}+\mathrm{EM}\) denote its classical extension by the law of excluded middle. Here excluded middle is expressed schematically as
\[
\varphi\lor\neg\varphi
\]
for propositions \(\varphi\) in the intended formulation; the precise language and scope of this schema must be specified by that formulation.

\emph{Problem.} Determine the exact proof-theoretic strength of each of the theories
\[
\mathrm{MF}\qquad\text{and}\qquad\mathrm{MF}+\mathrm{EM}.
\]

\emph{Unresolved specifications.} The supplied record does not give the axioms, inference rules, or precise presentation of \(\mathrm{MF}\), nor does it specify how excluded middle is added in that presentation. It also does not define the intended measure of proof-theoretic strength. Its accompanying assessment mentions an exact ordinal or proof-theoretic calibration, without specifying an ordinal invariant, reference theories, or comparison relation. Consequently, the question above preserves the supplied target, but a fully specified formal problem requires these choices to be recovered from the cited sources; no particular choices or equivalence between possible calibrations are asserted here.
\subsection{Short English statement}
Determine exactly how strong the Minimalist Foundation and its extension by excluded middle are in proof-theoretic terms. The supplied record leaves the formal presentation and intended measure of strength unspecified. Its open status is supplied dated curation, not an independently verified conclusion.
\subsection{Sources}
\begin{itemize}
\item Local UnsolvedMath registry: supplied record with permanent ID \(30004553\), title \emph{Proof-Theoretic Strength of Minimalist Foundations}, in the Oberwolfach Reports Open Problems collection. Exact source problem code: \texttt{catalogue entry}. The record gives proposed year 2020 and supplies no individual proposer attribution.
\item Supplied source citation: \emph{Mathematical Logic: Proof Theory, Constructive Mathematics} (2021), Oberwolfach Reports, pp.\ 1693--1757, DOI \texttt{10.4171/owr/2020/34}. \url{https://doi.org/10.4171/owr/2020/34}. The background describes Oberwolfach Reports as documenting mathematical workshops and their open-problem discussions.
\item Supplied literature reference: Maietti--Sabelli, 2024. \url{https://doi.org/10.1016/j.apal.2024.103524}. No article title or relevant theorem numbers are supplied.
\item Supplied literature reference: Oberwolfach Report 34/2020. \url{https://doi.org/10.4171/owr/2020/34}.
\item Literature assessment dated \texttt{2026-08-22}: the supplied curation labels the problem open and reports that recent work establishes equiconsistency and useful upper bounds without identifying the exact ordinal or proof-theoretic calibration for \(\mathrm{MF}\) and \(\mathrm{MF}+\mathrm{EM}\). This is supplied dated curation derived from the supplied OWR curation, not independently verified source inspection or a new literature assessment. The record explicitly calls for independent verification of time-sensitive status and cited claims.
\end{itemize}
\end{document}
